%% CV for Eltwin Group - Embedded Software Engineer (Dansk)
%% Compiled with: lualatex

\documentclass[11pt,a4paper,sans]{moderncv}
\moderncvstyle{banking}
\moderncvcolor{blue}

\renewcommand*{\firstnamestyle}[1]{{\fontsize{34}{36}\bfseries\upshape\color{color1}#1}}
\renewcommand*{\lastnamestyle}[1]{{\fontsize{34}{36}\bfseries\upshape\color{color1}#1}}
\renewcommand*{\sectionstyle}[1]{{\sectionfont\color{color1}#1}}

\usepackage[utf8]{inputenc}
\usepackage{hyperref}
\hypersetup{
    colorlinks=true,
    linkcolor=blue,
    filecolor=magenta,
    urlcolor=blue,
    pdftitle={Jacob El-Omar - CV},
    pdfpagemode=FullScreen,
}
\usepackage[scale=0.80]{geometry}
\usepackage{import}

\name{Jacob}{El-Omar}
\address{Aalborg, Danmark}{}{}
\phone[mobile]{+45 53 35 27 82}
\email{elomarjc@gmail.com}
\extrainfo{\href{https://linkedin.com/in/jacob-el-omar}{LinkedIn} \textbullet{} \href{https://github.com/elomarjc}{GitHub}}

\begin{document}

\makecvtitle

% ============================================================
%     PROFILE
% ============================================================

\vspace{6pt}
\small{Nyuddannet civilingeniør (MSc) i Elektroniske Systemer fra Aalborg Universitet med speciale i reeltids-embedded software, motorstyrings-algoritmer og firmware til mikrocontrollere. Erfaring med C/C++, data-pipelines i Python, hardware-software integration og fejlfinding. Kompetencer i at anvende agentiske AI-workflows med Claude Code til hurtig testautomatisering, refakturering og teknisk dokumentation. Motiveret for at bidrage til udviklingen af energieffektive elektroniske styresystemer hos Eltwin Group i Risskov.}

% ============================================================
%     TECHNICAL SKILLS
% ============================================================

\section{Faglige Kompetencer}
\vspace{1pt}
\begin{itemize}

\item \textbf{Embedded Software \& Programmering}: Embedded C/C++, Python (NumPy, SciPy), C\#, Cypress PSoC mikrocontrollere, I2C/SPI/UART kommunikation, reeltidsfejlfinding.

\item \textbf{Reguleringsteknik \& Motorstyring}: Motorstyrings-algoritmer, tilstandsestimering (Kalman-filtrering), LQR/PID reguleringssløjfer samt systemmodellering i MATLAB \& Simulink.

\item \textbf{Softwarekvalitet \& Metoder}: Versionsstyring med Git, automatiserede test-pipelines, hardware-in-the-loop (HIL) testmetoder og teknisk dokumentation.

\item \textbf{AI-Værktøjer \& PBL-samarbejde}: Anvendelse af Claude Code til kodegenerering og testautomatisering. Erfaring med tværfagligt projektsamarbejde via problembaseret læring (PBL).

\end{itemize}

% ============================================================
%     EDUCATION
% ============================================================

\section{Uddannelse}
\vspace{1pt}

\cventry{Sep 2023--Jun 2025}{MSC i Elektroniske Systemer}{Aalborg Universitet}{Aalborg, Danmark}{}{
Kandidatspeciale: ``Adaptive Noise Cancellation For Electronic Stethoscopes.'' Modellering af komplekse stokastiske signaler, reeltids-signalbehandling, adaptiv filtrering og algoritmeudvikling i Python og MATLAB.
}

\vspace{3pt}

\cventry{Sep 2020--Aug 2023}{BSc i Elektronik og IT}{Aalborg Universitet}{Aalborg, Danmark}{}{
Specialisering inden for reguleringsteknik, embedded mikrocontrollere og signalbehandling.
}

% ============================================================
%     PROFESSIONAL EXPERIENCE
% ============================================================

\section{Erhvervserfaring}
\vspace{3pt}

\cventry{Sep 2024--Jun 2025}{AI/ML Engineer (Praktikant)}{Ai Health Highway}{Aalborg, Danmark / Remote}{}{
\begin{itemize}
    \item Udviklede og validerede signalbehandlings- og dataanalysealgoritmer i Python til støjreduktion i reeltid.
    \item Håndterede og rensede store akustiske datasæt i overensstemmelse med medicinske compliance-standarder.
    \item Udarbejdede teknisk dokumentation og testafrapporteringer til tværfaglige teams.
\end{itemize}}

\vspace{3pt}

\cventry{Jan 2026--Nuværende}{IT- \& Webudvikler (Frivillig)}{Dawah Danmark}{Aalborg, Danmark}{}{
\begin{itemize}
    \item Moderniserede, strukturerede og vedligeholdt databaser og webinfrastruktur (WordPress, MySQL).
    \item Forbedrede dataopsamling og systemarkitektur for at optimere de interne arbejdsgange.
\end{itemize}}

\vspace{3pt}

\cventry{Jun 2023--Aug 2024}{Lager- og Logistikmedarbejder}{Lyn Express}{Aalborg, Danmark}{}{
\begin{itemize}
    \item Håndterede registrering, sortering og overvågning af logistikdata under højt arbejdstempo.
\end{itemize}}

% ============================================================
%     SELECTED PROJECTS
% ============================================================

\section{Udvalgte Projekter}
\vspace{1pt}
\begin{itemize}

\item{\textbf{Grow a Fish (Fritidsprojekt)}: Egenudviklet full-stack mobilspil med vækstdynamikker, multiplayer og in-game butik (Flutter, Dart, SQL, REST API).}

\item{\textbf{AAUSAT6 ADCS Development}: Udvikling af attituderegulatorer til CubeSat-stabilisering ved hjælp af LQR og robust regulering, herunder B-dot-algoritmer (MATLAB, Simulink, LQR, Kalman-filter).}

\item{\textbf{CoM Calibration of 3-DoF Satellite Simulator}: Algoritmeudvikling til estimering af tyngdepunkt (CoM) ved brug af kinematik og Kalman-filtrering (MATLAB, Kalman-filter, Kinematik, OnShape, CNC).}

\item{\textbf{Robot Collision-Avoidance}: Prædiktivt navigationssystem til MiR200-platforme ved brug af UWB og regression (UWB, Kalman-filter, Regression, MiR200).}

\end{itemize}

% ============================================================
%     LANGUAGES
% ============================================================

\section{Sprog}
\vspace{1pt}
\begin{itemize}
\item Dansk (Modersmål), Engelsk (Flydende i skrift og tale).
\end{itemize}

% ============================================================
%     REFERENCES
% ============================================================

\section{Referencer}
\vspace{1pt}
\begin{itemize}
\item Kan oplyses efter anmodning.
\end{itemize}

\end{document}
